\section{Dataset and Preprocessing}
\label{sec:data}
% DRAFT (2026-06-17). Numbers verified from results/ and notebook 01_eda. Own the prose.

\textbf{Dataset.}
We use DGA records from transformer main tanks: 4{,}563 samples from 628 transformers and 42
manufacturers, sampled between 2019 and 2024 (median six samples per unit over $\sim$4.5 years).
Rated voltages are recorded as free-text multi-winding strings (e.g.\ ``115/22'') and are not
used as model inputs. A per-gas summary is given in Table~\ref{tab:gases}.

\textbf{Data provenance.}
The dataset was assembled for this study from operational main-tank DGA records of in-service
transformers in Thailand (multiple operational sources). Unit identifiers were anonymised, and
no nameplate or location information is used. The raw data are not publicly available; an
anonymised extract (hashed unit identifiers, dates, gas concentrations, and binary event flags)
is released with our code for reproducibility.

\begin{table}[t]
\centering
\caption{Per-gas summary over the 4{,}563 samples (concentrations in ppm).}
\label{tab:gases}
\begin{tabular}{lrrr}
\toprule
Gas & Median & Max & \% zero \\
\midrule
H2   & 16    & 2{,}537  & 1.1 \\
CH4  & 18    & 2{,}469  & 10.1 \\
C2H2 & 0     & 954      & 97.4 \\
C2H4 & 2     & 3{,}394  & 39.1 \\
C2H6 & 14    & 1{,}085  & 9.5 \\
CO   & 193   & 2{,}437  & 0.8 \\
CO2  & 2{,}281 & 23{,}574 & 0.7 \\
\bottomrule
\end{tabular}
\end{table}

\textbf{Features.}
Each sample provides seven diagnostic gases used as model features: H2, CH4, C2H2, C2H4,
C2H6, CO and CO2. Atmospheric oxygen and nitrogen, and the total-combustible-gas sum, are
excluded to avoid redundancy and leakage.

\textbf{Data characteristics.}
After cleaning, the gases have essentially no missing values; sparsity instead appears as
zeros---acetylene is zero in $97.4\%$ of samples, consistent with arcing being rare. The
concentrations are heavily right-skewed and remain skewed for C2H2 and C2H4 even after a
logarithmic transform. One oxygen reading is a clear data-entry error (about
$7.2\times10^{8}$~ppm), which motivates outlier clipping.

\textbf{Preprocessing.}
The pipeline is configuration-driven: each gas is clipped at its 99.9th percentile, missing
readings are imputed as zero (taken as below the detection limit), a $\log(1+x)$ transform
is applied, and the features are standardised.

\textbf{Conventional labels (weak, external).}
For comparison we compute Duval Triangle~1, IEC~60599 and Rogers labels with our reference
implementation. The Duval class distribution (over all samples) is imbalanced: PD~1511, T1~1451,
T2~686, T3~492, D2~30, D1~24, DT~16, with about $7.7\%$ not determined. Many of the PD samples are
methane-dominant cases that Triangle~1, lacking a stray-gassing zone, labels PD (the Duval Pentagon
adds such a zone); this is mostly a per-sample effect -- at the per-unit (latest-state) level only
about 100 units are PD. We therefore treat the conventional labels as indicative, not ground truth.

\textbf{Validation notes.}
391 of 4{,}563 samples ($8.6\%$) carry a raw free-text field-event note (for example a Buchholz
trip or a bushing failure). Filtered to genuine protection/damage events and aggregated to units,
these give the weak unit-level label \texttt{fault\_note} used for validation (Section~\ref{sec:exp}); the
sample-level $8.6\%$ and the unit-level label rate are therefore different quantities. Most
notes are partly in Thai ($69\%$; $127$ are Thai-only) and invisible to an English keyword
filter; a curated translation of all 211 distinct Thai notes (released with the code, pending
domain validation) identifies 15 additional event notes, whose effect is evaluated in
Section~\ref{sec:results}.
